
\section{Resumo e verificação}
\label{sec:dk_resumo}

Combinando os resultados obtidos nos planos lateral e horizontal, é possível determinar a posição tridimensional do centro de preensão do myCobot~280 em função dos ângulos articulares e dos parâmetros geométricos do manipulador.

Considerando $D=L_6+L_7$ e os ângulos absolutos definidos na Equação~\eqref{eq:angulos}, a formulação da cinemática direta é dada por:

\begin{equation}
    \boxed{
    \begin{aligned}
        d &= L_2\cos\theta_A
        + L_3\cos\theta_B
        + L_4\cos\theta_C
        + D\cos q_5\cos\alpha\\
        h &= L_1
        + L_2\sin\theta_A
        + L_3\sin\theta_B
        + L_4\sin\theta_C
        + D\cos q_5\sin\alpha\\
        l &= L_5-D\sin q_5\\[4pt]
        x &= l\cos q_1-d\sin q_1\\
        y &= l\sin q_1+d\cos q_1\\
        z &= h
    \end{aligned}
    }
    \label{eq:dk_final}
\end{equation}

Estas expressões permitem calcular a posição do centro de preensão relativamente ao referencial da base, a partir dos ângulos das juntas e dos comprimentos dos elos.

Adicionalmente, a direção longitudinal da garra pode ser determinada através das componentes $-\sin q_5$, $\cos q_5\cos\alpha$ e $\cos q_5\sin\alpha$, expressas no referencial associado ao plano do braço. Estas componentes são posteriormente transformadas para o referencial da base, considerando a rotação $q_1$.

\subsection{Verificação dos resultados}

Para verificar o comportamento do modelo matemático, a cinemática direta foi avaliada em duas configurações distintas do manipulador.

\subsubsection{Configuração de referência}

A primeira configuração corresponde à posição inicial do manipulador, designada por \textit{pose zero}, na qual todas as juntas apresentam um ângulo de $0^\circ$.

De acordo com a Equação~\eqref{eq:angulos}, nesta configuração obtêm-se:

\begin{equation}
    \theta_A=\theta_B=\theta_C=90^\circ,
    \qquad \alpha=0^\circ
\end{equation}

Assim, o braço encontra-se orientado verticalmente, enquanto a garra permanece na horizontal, apontando no sentido positivo do eixo $y$.

As expressões da cinemática direta simplificam-se para:

\begin{equation}
    d=L_6+L_7, \qquad
    h=L_1+L_2+L_3+L_4
    \qquad \Rightarrow \qquad
    (x,y,z)=(L_5,d,h)
    \label{eq:pose_zero}
\end{equation}

Substituindo os valores dos comprimentos apresentados na Tabela~\ref{tab:elos}, obtém-se:

\begin{equation}
    (x,y,z)=(0{,}0646;\ 0{,}1513;\ 0{,}4111)
    \ \text{m}
\end{equation}

A direção longitudinal da garra é representada pelo vetor $(0,1,0)$, estando o resultado de acordo com a geometria definida para a configuração de referência.

\subsubsection{Configuração de trabalho}

A segunda configuração corresponde a uma posição de trabalho, na qual a garra se encontra orientada verticalmente para baixo.

Os ângulos articulares considerados são:

\begin{equation}
    \mathbf{q}=
    (18{,}84^\circ;\
    -12{,}74^\circ;\
    -94{,}21^\circ;\
    16{,}95^\circ;\
    0^\circ;\
    0^\circ)
    \label{eq:q_trabalho}
\end{equation}

Aplicando as relações definidas para os ângulos absolutos, obtêm-se:

\begin{equation}
    \theta_A=77{,}26^\circ,\qquad
    \theta_B=-16{,}95^\circ,\qquad
    \theta_C=0^\circ,\qquad
    \alpha=-90^\circ
\end{equation}

Substituindo estes valores nas equações do plano lateral, obtêm-se as coordenadas:

\begin{equation}
    d=0{,}1893\ \text{m},\qquad
    h=0{,}0600\ \text{m}
\end{equation}

Considerando a rotação da base e o desvio lateral do manipulador, a posição tridimensional do centro de preensão é:

\begin{equation}
    (x,y,z)\approx
    (0{,}0000;\ 0{,}2000;\ 0{,}0600)
    \ \text{m}
\end{equation}

Deste modo, o centro de preensão encontra-se aproximadamente a $20$ cm de distância horizontal da base, na direção positiva do eixo $y$, e a uma altura de $6$ cm.

A Figura~\ref{fig:poses_fk} apresenta as duas configurações consideradas na verificação da cinemática direta.

\begin{figure}[H]
    \centering
    \resizebox{0.75\textwidth}{!}{
        % Pose zero e configuração de trabalho (plano do braço); unidades em cm
\begin{tikzpicture}[scale=0.26, >=stealth,
    elo/.style={line width=2pt, line cap=round},
    junta/.style={circle, draw, fill=white, line width=0.9pt, inner sep=1.8pt},
    cota/.style={gray!70!black, thin},
    ang/.style={->, thin},
    every node/.style={font=\footnotesize}]

  \draw[->, thin] (-3,0) -- (22,0) node[below left] {$d$};
  \draw[->, thin] (0,0) -- (0,46) node[left] {$h$};
  \fill[gray!25] (-2.5,-0.8) rectangle (2.5,0.0);
  \draw[elo, blue!70!black] (0.0,0.0) -- (0.0,13.2) -- (0.0,24.2) -- (0.0,33.8) -- (0.0,41.1) -- (15.1,41.1);
  \node[junta] at (0.0,13.2) {};
  \node[junta] at (0.0,24.2) {};
  \node[junta] at (0.0,33.8) {};
  \node[junta] at (0.0,41.1) {};
  \fill (15.1,41.1) circle (5pt);
  \node[right=4pt, align=left] at (13.1,37.1) {garra a apontar para $+y$\\$(x, y, z) = (0{,}0646;\ 0{,}1513;\ 0{,}4111)$ m};
  \node[font=\small\bfseries] at (8,50) {Pose zero};
  \draw[->, thin] (35,0) -- (62,0) node[below left] {$d$};
  \draw[->, thin] (38,0) -- (38,46) node[left] {$h$};
  \fill[gray!25] (35.5,-0.8) rectangle (40.5,0.0);
  \draw[elo, orange!85!black] (38.0,0.0) -- (38.0,13.2) -- (40.4,23.9) -- (49.6,21.1) -- (56.9,21.1) -- (56.9,6.0);
  \node[junta] at (38.0,13.2) {};
  \node[junta] at (40.4,23.9) {};
  \node[junta] at (49.6,21.1) {};
  \node[junta] at (56.9,21.1) {};
  \fill (56.9,6.0) circle (5pt);
  \node[below=4pt, align=center] at (56.9,6.0) {garra a apontar para baixo\\$(x, y, z) = (0;\ 0{,}2000;\ 0{,}0600)$ m};
  \node[font=\small\bfseries] at (50,50) {Configuração de trabalho};
\end{tikzpicture}
    }
    \caption{Configuração de referência e configuração de trabalho utilizadas na verificação da cinemática direta.}
    \label{fig:poses_fk}
\end{figure}

\subsubsection{Validação no CoppeliaSim}

A validação do modelo cinemático é realizada através do modo de diagnóstico implementado em MATLAB, comparando a posição do centro de preensão calculada pela cinemática direta com a posição obtida diretamente no CoppeliaSim.

O erro de posicionamento é determinado através da distância euclidiana entre as duas posições, sendo expresso em milímetros:

\begin{equation}
    e_{\mathrm{modelo}} =
    1000
    \left\|
    \mathbf{p}_{\mathrm{DK}} -
    \mathbf{p}_{\mathrm{sim}}
    \right\|_2
    \label{eq:erro_modelo}
\end{equation}

onde $\mathbf{p}_{\mathrm{DK}}$ representa a posição calculada pelo modelo de cinemática direta e $\mathbf{p}_{\mathrm{sim}}$ corresponde à posição do centro de preensão obtida no CoppeliaSim, ambas expressas em metros e relativamente ao referencial da base.

No modo de diagnóstico, o programa obtém os ângulos articulares do manipulador, calcula a posição através da função de cinemática direta e compara o resultado com a posição fornecida pelo simulador.

Nos resultados reportados para as configurações testadas, as diferenças observadas situaram-se entre $1{,}2$ mm e $4{,}0$ mm, indicando uma boa aproximação entre o modelo geométrico desenvolvido e o comportamento do manipulador em simulação.

Adicionalmente, o programa permite calcular o erro entre a posição pretendida e a posição efetivamente alcançada, possibilitando avaliar o desempenho do sistema de posicionamento.

Foi ainda considerada a consistência entre a cinemática direta e a cinemática inversa, verificando se a aplicação sucessiva dos dois modelos permite recuperar a posição inicialmente pretendida.

\clearpage
